\section{Conclusion}
\label{sec:conclusion}

We tested whether an explicit, secret-blind outline reduces involuntary story leakage and whether a carrier-defined residual direction mediates the effect. \gemma reproduces leakage when story content is unconstrained: the open anchor yields a $+16.33$ point order-balanced influence difference. A neutral brief then removes the detectable signal before an outline is introduced. Adding the outline does not reliably improve on either the brief alone or length-matched irrelevant context. Plot-twist evidence is underpowered, residual gaps are inconsistent, and projection removal destroys basic generation utility.

The central takeaway is that content-choice freedom, not demonstrated planning offload, explains the strongest contrast in this experiment. Future work should first screen broad premises that retain 65--80\% anchor discrimination, then use at least 100 pairs per condition, two calibrated judge families, and blinded human validation. Larger plot sets should be calibrated for premise-only guessability. Mechanistic studies should learn subspaces on disjoint secrets, compare multiple removal methods such as LEACE or INLP, and reject any intervention that materially changes length, coherence, or task completion.

